\uuid{nR2F}
\titre{Théorème de Bézout}
\auteur{bodin ; croizat ; sacre}
\organisation{exo7}
\licence{CC BY-NC-SA 4.0}
\datecreate{2022-03-08}
\niveau{L1}
\module{Algèbre}
\chapitre{Arithmétique dans Z}
\sousChapitre{Nombres premiers, nombres premiers entre eux}
\difficulte{3}
\theme{Théorème de Bézout}
\exo7qcmid{hilisit.arithmetique.025}

\contenu{
\begin{question}
\qclassification{103.04}
 Soit $n$ un entier tel que $5n$ soit un multiple de $7$. Quelles sont alors les affirmations vraies ?
    \begin{answers}
        \good{$n$ est un multiple de $7$.}
        \bad{$5$ divise $7n$.}
        \good{$7$ divise $n$.}
        \bad{$35$ divise $n$.}
    \end{answers}
    \begin{explanations}
    D'après le lemme de Gauss, puisque $7 | 5n$ et que $\pgcd(5,7) = 1$, on a $ 7 | n $ : ceci revient à dire que $n$ est un multiple de $7$.\\
    En revanche si l'on prend $n=7$, on constate que $5n = 35$ est bien multiple de $7$ mais que $5$ ne divise pas $7n=49$ et que $35$ ne divise pas $n=7$.
    \end{explanations}
\end{question}
}
